\chapter{الفضاءات المتجهية}\label{ch:b1:vspaces}

يبدأ الجبر الخطي هنا: فبديهيات الفضاءات المتجهية تعزل ما
يشترك فيه $\R^2$ و $\R^3$ وفضاءات كثيرات الحدود وفضاءات الدوال
--- إذ يمكن الجمع والضرب في سلّم. ويبني فصلان النظرية
(\cref{ch:b1:findim} يضيف البُعد)؛ واللغة التي يهيّئانها ---
\omterm{def:b1:vspaces:span}{الفضاء المولَّد}، \omterm{def:b1:vspaces:free}{والعائلة الحرة}، \omterm{def:b1:vspaces:free}{والأساس}، \omterm{def:b1:vspaces:sum}{والمجموع المباشر} --- هي الخبز اليومي لكل
فصل بعدهما. وفي كل ما يأتي، يرمز $K$ إلى $\R$ أو $\C$
(\emph{السلالم}).

\section{التعريف والأمثلة}

\begin{definition}[فضاء متجهي]\label{def:b1:vspaces:def}
\emph{الفضاء المتجهي على $K$}\index{فضاء متجهي} \omterm{def:b1:logic:sets}{مجموعة} $E$ مزوَّدة
بجمع يجعل $(E, +)$ زمرةً \omterm{def:b1:structures:group}{تبديلية} (وتُكتب صفرها $0_E$ أو
$0$)، وبضرب سلّمي $K \times E \to E$ بحيث يكون، من أجل
كل $\lambda, \mu \in K$ و $x, y \in E$:
\[
\lambda(x + y) = \lambda x + \lambda y,\quad
(\lambda + \mu) x = \lambda x + \mu x,\quad
\lambda(\mu x) = (\lambda\mu) x,\quad
1\,x = x .
\]
والنتائج: $0\,x = 0_E$ و $\lambda\,0_E = 0_E$ و $(-1)x = -x$، و
$\lambda x = 0_E \implies \lambda = 0$ أو $x = 0_E$ (بالضرب في
$\lambda^{-1}$).
\end{definition}

\begin{proof}[برهان النتائج]
من أجل $0\,x = 0_E$: من $(0 + 0)x = 0x + 0x$ و $(0+0)x = 0x$،
اختصر $0x$ في \omterm{def:b1:structures:group}{الزمرة} $(E, +)$. ومن أجل $\lambda\,0_E$: الحيلة
نفسها على $\lambda(0_E + 0_E)$. ومن أجل $(-1)x$: أضف $x$،
\[
x + (-1)x = 1\,x + (-1)x = \bigl(1 + (-1)\bigr)x = 0\,x = 0_E ,
\]
ومنه فالمقدار $(-1)x$ هو النظير الجمعي للمقدار $x$. وأخيرًا إذا كان $\lambda x
= 0_E$ مع $\lambda \neq 0$: فاضرب في $\lambda^{-1}$ (لأن
السلالم تكوّن \omterm{def:b1:structures:field}{حقلًا}) واستعمل البديهيتين
$\lambda^{-1}(\lambda x) = (\lambda^{-1}\lambda)x = 1x = x$
مع $\lambda^{-1}0_E = 0_E$: فيكون $x = 0_E$. وهذه القواعد الأربع، على
صغرها، تُستعمل صامتة في كل صفحة تأتي
--- والأخيرة هي بالضبط حيث تلزم الحقول:
فعلى السلالم $\Z$، سينتهك «الفضاء» $\Z/2\Z$ ذلك
بالمقدار $2\,x = 0$.
\end{proof}

\begin{example}\label{ex:b1:vspaces:examples}
$K^n$ (بالعمليات إحداثيةً إحداثية)؛ وكثيرات الحدود $K[X]$؛ و
الدوال $\mathcal{F}(A, K)$ من أيّ \omterm{def:b1:logic:sets}{مجموعة} $A$ إلى $K$ (بالعمليات
نقطةً نقطة) --- وهي تحتوي الدوال \omterm{def:b1:continuity:continuous}{المتصلة} والمتتاليات
$\mathcal{F}(\N, \R)$ وغيرها؛ و $\C$ فضاءً متجهيًا على $\R$. وفي كل
حالة تُورَّث البديهيات من بديهيات $K$.
\end{example}

\begin{definition}[الفضاء الجزئي]\label{def:b1:vspaces:subspace}
تكون $F \subseteq E$ \emph{فضاءً جزئيًا}\index{فضاء جزئي} إذا كان $0_E \in
F$ وكانت $F$ مستقرة بالجمع وبالضرب السلّمي ---
وبكيفية مكافئة:
\[
F \neq \emptyset
\qquad\text{و}\qquad
\forall x, y \in F,\ \forall \lambda \in K,\quad
x + \lambda y \in F .
\]
والفضاء الجزئي هو نفسه \omterm{def:b1:vspaces:def}{فضاء متجهي}. وكل تقاطع لفضاءات جزئية
فضاء جزئي؛ وأمّا الاتحاد فلا يكون كذلك في الغالب الأعمّ (بالبرهان نفسه في
\cref{exo:b1:structures:6}).
\end{definition}

\begin{example}\label{ex:b1:vspaces:subspaces}
في $\mathcal{F}(\R, \R)$: الدوال \omterm{def:b1:continuity:continuous}{المتصلة}، و
الدوال \omterm{def:b1:derivative:def}{القابلة للاشتقاق}، وكثيرات الحدود من الدرجة $\leq n$ (وتُكتب
$K_n[X]$ داخل $K[X]$)، وحلول معادلة تفاضلية خطية
متجانسة (وهذا بالضبط ما قاله \cref{thm:b1:diffeq:homogeneous2}).
وأمثلة مضادة: $\{f : f(0) = 1\}$ (فلا صفر فيها)، والدرجة $n$
بالضبط (فهي غير مستقرة بالجمع).
\end{example}

\begin{example}[فضاء جزئي أو لا: أربعة أحكام، محاجَجة]
\label{ex:b1:vspaces:subspaceornot}
في فضاء المتتاليات الحقيقية:
\begin{itemize}
  \item $\{u : u \text{ محدودة}\}$ \emph{هي} \omterm{def:b1:vspaces:subspace}{فضاء جزئي}: فالمقدار $0$
        محدود، وإذا كان $\abs{u_n} \leq M$ و $\abs{v_n} \leq M'$،
        فإن $\abs{u_n + \lambda v_n} \leq M + \abs\lambda M'$.
  \item $\{u : u_n \to 1\}$ \emph{ليست} كذلك: فالمتتالية المعدومة
        مفقودة (ومجموع عضوين يؤول إلى $2$).
  \item $\{u : u \text{ رتيبة}\}$ \emph{ليست} كذلك: فالمتتاليتان $u_n = n$
        و $v_n = -n + (-1)^n$ رتيبتان، ومجموعهما
        $(-1)^n$ ليس كذلك؛ فالاستقرار بالجمع هو البديهية
        التي تفشل، وإن كانت \omterm{def:b1:logic:sets}{المجموعة} تحتوي $0$ وكل
        مضاعفات أعضائها السلّمية.
  \item $\{u : u_{n+1} = u_n^2\}$ \emph{ليست} كذلك: فهي تحتوي $0$
        لكن $2u$ يفلت ما إن يكون $u$ عضوًا غير معدوم
        ($2u_{n+1} \neq (2u_n)^2$ عمومًا) --- فالتربيع هو
        اللاخطية.
\end{itemize}
وترتيب العمل هو نفسه دائمًا: اختبر $0$ أولًا (فهو الأرخص)،
ثم الاستقرار --- وللنقض، يغلب زوجٌ مضاد صريح واحد
أيَّ قدر من الشكّ.
\end{example}

\section{الفضاء المولَّد والمجاميع والمجاميع المباشرة}

\begin{definition}[التركيبات الخطية والفضاء المولَّد]\label{def:b1:vspaces:span}
\emph{التركيبة الخطية} للعائلة $(x_1, \dots, x_p)$ من
متجهات $E$ هي أيّ $\lambda_1 x_1 + \dots + \lambda_p x_p$
($\lambda_i \in K$). ومجموعتها كلها هي
\emph{الفضاء المولَّد}\index{فضاء مولَّد} $\operatorname{Vect}(x_1, \dots, x_p)$: وهو
\omterm{def:b1:vspaces:subspace}{فضاء جزئي}، وأصغرُ \omterm{def:b1:vspaces:subspace}{فضاء جزئي} يحتوي العائلة.
\end{definition}

\begin{proof}[برهان الادعاءين]
الاستقرار: مجموع تركيبتين خطيتين $\sum\lambda_i x_i +
\sum\mu_i x_i = \sum(\lambda_i + \mu_i)x_i$ هو مرة أخرى تركيبة خطية، وكذلك
المضاعف السلّمي $\mu\sum\lambda_i x_i = \sum(\mu\lambda_i)
x_i$؛ وتبيّن التركيبة المعدومة أن $0$ ينتمي: ومنه \omterm{def:b1:vspaces:span}{فالفضاء المولَّد}
\omterm{def:b1:vspaces:subspace}{فضاء جزئي}. والأصغرية: ليكن $H$ أيَّ \omterm{def:b1:vspaces:subspace}{فضاء جزئي} يحتوي $x_1,
\dots, x_p$. فبالاستقرار بالضرب السلّمي يكون كل
$\lambda_i x_i \in H$، وبالاستقرار بالجمع يقع مجموعها
في $H$: ومنه فكل تركيبة خطية تنتمي إلى $H$، أي
$\operatorname{Vect}(x_1, \dots, x_p) \subseteq H$. ومنه \omterm{def:b1:vspaces:span}{فالفضاء المولَّد}
محتوى في كل \omterm{def:b1:vspaces:subspace}{فضاء جزئي} يحتوي العائلة: فهو
الأصغر.
\end{proof}

\begin{definition}[المجموع والمجموع المباشر]\label{def:b1:vspaces:sum}
من أجل فضاءين جزئيين $F, G$ من $E$:
\[
F + G = \{\,u + v : u \in F,\ v \in G\,\}
\]
هو \omterm{def:b1:vspaces:subspace}{فضاء جزئي} (وهو الأصغر الذي يحتوي $F \cup G$). ويكون المجموع
\emph{مباشرًا}\index{مجموع مباشر}، ويُكتب $F \oplus G$، إذا تفكّك كل
عنصر من $F + G$ تفكّكًا \emph{وحيدًا} في صورة $u + v$؛
وبكيفية مكافئة (انظر أدناه) إذا كان $F \cap G = \{0\}$. وإذا كان $E = F \oplus
G$، يكون الفضاءان الجزئيان \emph{متتامّين}\index{فضاءان
متتامّان} في $E$.
\end{definition}

\begin{example}[مجموع مستقيمين]
\label{ex:b1:vspaces:sumoflines}
في $\R^3$، لتكن $F = \operatorname{Vect}\bigl((1,0,1)\bigr)$ و
$G = \operatorname{Vect}\bigl((0,1,1)\bigr)$. مجموعهما هو
\[
F + G = \{\,a(1,0,1) + b(0,1,1)\,\}
= \{(a,\ b,\ a + b)\} = \{(x, y, z) : z = x + y\},
\]
أي المستوي المارّ بالمبدأ والحاوي للمستقيمين. وهو
أكبر قطعًا من الاتحاد $F \cup G$ (وهو مجرد تقاطع
المستقيمين): فالمتجهة $(1, 1, 2) = (1,0,1) + (0,1,1)$ تقع في
المجموع ولا تقع على أيّ من المستقيمين. و $F \cap G = \{0\}$ (لأن
متجهةً مشتركة تقتضي $a(1,0,1) = b(0,1,1)$، وتفرض إحداثيتاها
الأوليان $a = b = 0$): ومنه \omterm{def:b1:vspaces:sum}{فالمجموع مباشر}، و $F \oplus
G$ هو ذلك المستوي بالضبط.
\end{example}

\begin{proposition}\label{prop:b1:vspaces:directsum}
يكون $F + G$ مباشرًا إذا وفقط إذا كان $F \cap G = \{0\}$.
\end{proposition}

\begin{proof}
إذا وقع $w \neq 0$ ما في $F \cap G$: فإن $w = w + 0 = 0 + w$ تفكيكان
للمقدار $w$. وبالعكس، إذا كان $u + v = u' + v'$ مع $u, u'
\in F$ و $v, v' \in G$، فإن $u - u' = v' - v$ ينتمي إلى $F \cap G =
\{0\}$: ومنه فالتفكيكات وحيدة.
\end{proof}

\begin{method}[البرهان على $E = F \oplus G$]
\label{met:b1:vspaces:directsum}
شيئان للتحقق، ولكلٍّ حركته الافتتاحية المعيارية.
\begin{enumerate}
  \item \emph{التقاطع البديهي.} خذ $x \in F \cap G$،
        واكتب شرطَي الانتماء صراحةً، واحصر $x =
        0$. (ولا تحاجج بصورة أبدًا: قارن المزالق أدناه.)
  \item \emph{المجموع هو كل شيء.} خذ $x \in E$ كيفيًا
        و\emph{أنتج} التفكيك $x = f + g$ ---
        إمّا بتخمين $f$ من الهدف (فعلى $f$ أن يحقق
        الخاصية المعرِّفة للمقدار $F$، وهي تملي صيغته
        عادةً) وإمّا بحلّ الجملة الخطية التي تعبّر عن $x$
        إزاء مولّدات $F$ و $G$.
\end{enumerate}
وحين تُخمَّن صيغة التفكيك، تكون الوحدانية
تلقائية من الخطوة 1؛ وحين لا يكون الوجود وحده واضحًا، ففي الخطوة 2
يقع العمل. ويجري المثالان أدناه هذه الطريقة: فمن أجل
الدوال الزوجية والفردية تفرض صيغةَ $f$ تقويمُ المتطابقة المرجوّة
عند $x$ و $-x$؛ ومن أجل كثيرات الحدود المنعدمة عند نقطة، التقويمُ عند $a$.
\end{method}

\begin{example}\label{ex:b1:vspaces:evenodd}
في $\mathcal{F}(\R, \R)$، تكون الدوال الزوجية $\mathcal{P}$ و
الدوال الفردية $\mathcal{I}$ متتامّتين: إذ تُكتب أيّ $f$
\[
f(x) = \underbrace{\frac{f(x) + f(-x)}{2}}_{\text{زوجية}}
+ \underbrace{\frac{f(x) - f(-x)}{2}}_{\text{فردية}},
\]
والدالة الزوجية والفردية معًا معدومة. (وتطبيقًا على $\exp$، هذا هو
الزوج $(\cosh, \sinh)$ في \cref{ch:b1:functions}.)
\end{example}

\begin{example}[{زوج متتامّ في $K_n[X]$}]
\label{ex:b1:vspaces:hyperplane}
ثبّت $a \in K$ وضع $F = \{P \in K_n[X] : P(a) = 0\}$ و $G =
\operatorname{Vect}(1)$ (وهي الثوابت). عندئذ $K_n[X] = F \oplus
G$. فعلًا يتألف $F \cap G$ من الثوابت المنعدمة عند $a$،
أي $\{0\}$؛ وكل $P$ يتفكّك
\[
P = \underbrace{\bigl(P - P(a)\bigr)}_{\in F}
+ \underbrace{P(a)}_{\in G} .
\]
والتفكيك جدير بالحفظ: فطرح القيمة عند
نقطة هو الكيفية المعيارية للإسقاط على «الدوال المنعدمة عند
$a$». ولاحظ أن $F$ \omterm{def:b1:vspaces:subspace}{فضاء جزئي} كبير و $G$ صغير؛ فليس على
الزوج المتتامّ أن يكون متوازنًا بأيّ معنى.
\end{example}

\begin{example}[الفضاء الجزئي المتتامّ ليس وحيدًا أبدًا]
\label{ex:b1:vspaces:manysupplements}
في $\R^2$، لتكن $F = \operatorname{Vect}\bigl((1,0)\bigr)$ (وهو
محور $x$). فكلٌّ من $G = \operatorname{Vect}\bigl((0,1)\bigr)$ و
$G' = \operatorname{Vect}\bigl((1,1)\bigr)$ متتامّ مع
$F$: إذ يقاطع كلٌّ منهما $F$ عند $0$ فقط، ويعطي كل زوج مجموعًا هو $\R^2$.
ويختلف تفكيكا المتجهة نفسها:
\[
(2,\ 1.5) = \underbrace{(2, 0)}_{\in F} +
\underbrace{(0, 1.5)}_{\in G}
= \underbrace{(0.5,\ 0)}_{\in F} +
\underbrace{(1.5,\ 1.5)}_{\in G'} .
\]
وفي الواقع \emph{كل} مستقيم غير $F$ نفسه هو
متتامّ مع $F$ في $\R^2$: فالمتتامّات كثيرة،
والحديث عن المتتامّ «الوحيد» بلا معنى حتى تفرز
بنيةٌ إضافية (جداء سلّمي، \cref{ch:b1:euclid})
واحدًا منها.
\end{example}

\begin{omfigure}
\begin{tikzpicture}[scale=1.5]
  \draw[->, thin] (-0.4,0) -- (2.9,0) node[below] {$x$};
  \draw[->, thin] (0,-0.4) -- (0,1.9) node[left] {$y$};
  \draw[omDef, very thick] (-0.3,0) -- (2.7,0)
    node[below, pos=0.95, font=\small] {$F$};
  \draw[omThm, very thick] (0,-0.3) -- (0,1.7)
    node[left, pos=0.95, font=\small] {$G$};
  \draw[omProp, very thick] (-0.25,-0.25) -- (1.75,1.75)
    node[right, pos=0.97, font=\small] {$G'$};
  \fill (2,1.5) circle (0.045) node[above right, font=\small]
    {$(2, 1.5)$};
  \draw[omThm, dashed] (2,1.5) -- (2,0);
  \draw[omDef, dashed] (2,1.5) -- (0,1.5);
  \draw[omProp, dashed] (2,1.5) -- (0.5,0);
  \fill (2,0) circle (0.035);
  \fill (0.5,0) circle (0.035);
\end{tikzpicture}

{\small تفكيكان للنقطة نفسها من $\R^2$ على امتداد $F$
(وهو محور $x$): مع المتتامّ $G$ (بهبوط عمودي) ومع
المتتامّ $G'$ (بهبوط مائل). وتختلف مركّبتا $F$:
فالإسقاط يتعلق باتجاه النزول.}
\end{omfigure}

\section{العائلات الحرة والعائلات المولِّدة والأسس}

\begin{definition}\label{def:b1:vspaces:free}
تكون العائلة $(x_1, \dots, x_p)$ من متجهات $E$:
\begin{itemize}
  \item \emph{مولِّدة} (للمقدار $E$) إذا كان
        $\operatorname{Vect}(x_1,\dots,x_p) = E$؛
  \item \emph{حرة}\index{عائلة حرة} (ومتجهاتها \emph{مستقلة
        خطيًا}) إذا كان
        \[
        \lambda_1 x_1 + \dots + \lambda_p x_p = 0
        \implies \lambda_1 = \dots = \lambda_p = 0 ;
        \]
        وإلا فهي \emph{مرتبطة}؛
  \item \emph{أساسًا}\index{أساس} إذا كانت حرة ومولِّدة.
\end{itemize}
\end{definition}

\begin{proposition}[الإحداثيات]\label{prop:b1:vspaces:coordinates}
تكون $(e_1, \dots, e_n)$ \omterm{def:b1:vspaces:free}{أساسًا} للمقدار $E$ إذا وفقط إذا كان كل $x \in E$
تركيبةً \emph{وحيدة} $x = \lambda_1 e_1 + \dots +
\lambda_n e_n$؛ وتكون السلالم $\lambda_i$ هي
\emph{إحداثيات}\index{إحداثيات} $x$ في \omterm{def:b1:vspaces:free}{الأساس}.
\end{proposition}

\begin{proof}
كونها مولِّدة $=$ وجود التفكيك. والوحدانية $=$
الحرية: إذ يختلف تفكيكان للمقدار $x$ نفسه بتركيبة
تساوي $0$؛ وتفرض الحرية انعدام كل معاملاتها --- وهي
فروق \omterm{prop:b1:vspaces:coordinates}{الإحداثيات}. وبالعكس، تعطي تركيبة معدومة غير بديهية
التفكيكين $0 = \sum \lambda_i
e_i = \sum 0\,e_i$.
\end{proof}

\begin{example}\label{ex:b1:vspaces:bases}
\emph{\omterm{def:b1:vspaces:free}{الأساس} القانوني} في $K^n$: $e_i = (0, \dots, 1, \dots, 0)$
($1$ في الخانة $i$). والوحيدات الحدّ $(1, X, X^2, \dots, X^n)$: أساسٌ
للمقدار $K_n[X]$ (والحرية: لأن تركيبة معدومة هي \omterm{def:b1:poly:def}{كثير الحدود} المعدوم، ومنه
تنعدم كل المعاملات، \cref{def:b1:poly:def}). وفي
$\C$ على $\R$: \omterm{def:b1:vspaces:free}{الأساس} $(1, \iu)$.
\end{example}

\begin{remark}[الإحداثيات عمل جماعي]
\label{rem:b1:vspaces:teameffort}
تتعلق الإحداثية الأولى للمقدار $x$ في \omterm{def:b1:vspaces:free}{أساس} $(e_1, \dots, e_n)$
بجميع \emph{كل} متجهات \omterm{def:b1:vspaces:free}{الأساس}، لا بالمقدار $e_1$ وحده. ففي
$\R^2$: للمتجهة $(3, 1)$ إحداثية أولى $3$ في
\omterm{def:b1:vspaces:free}{الأساس} القانوني، لكن إحداثيتها الأولى $2$ في \omterm{def:b1:vspaces:free}{الأساس}
$\bigl((1,0), (1,1)\bigr)$ --- حُلَّ $(3,1) = a(1,0) + b(1,1)$:
$b = 1$ و $a = 2$. فتغيير متجهة \omterm{def:b1:vspaces:free}{أساس} واحدة يعيد خلط
\emph{كل} إحداثية؛ وسوف يعلّب \cref{ch:b1:matrices} إعادة
الخلط هذه في مصفوفة تغيير \omterm{def:b1:vspaces:free}{الأساس}.
\end{remark}

\begin{example}[اختبار أساس مرشّح من البداية إلى النهاية]
\label{ex:b1:vspaces:testbasis}
هل $\mathcal{F} = (1 + X,\ 1 + X^2,\ X + X^2)$ \omterm{def:b1:vspaces:free}{أساس} للمقدار
$\R_2[X]$؟ اكتب $u_1, u_2, u_3$ لكثيرات الحدود الثلاثة.
\emph{الحرية}: تعطي تركيبة معدومة $a\,u_1 + b\,u_2 + c\,u_3 =
0$، معاملًا معاملًا،
\[
a + b = 0, \qquad a + c = 0, \qquad b + c = 0 ;
\]
وبطرح الأولين، $b = c$، ثم يعطي الثالث $2b =
0$: أي $a = b = c = 0$، فهي \omterm{def:b1:vspaces:free}{حرة}. \emph{والتوليد}: بدل حلّ
ثلاث جمل، لاحظ التركيبة المتناظرة
\[
u_1 + u_2 - u_3 = (1 + X) + (1 + X^2) - (X + X^2) = 2 ,
\]
ومنه $1 = \frac12(u_1 + u_2 - u_3)$؛ ثم
\[
X = u_1 - 1 = \tfrac12\bigl(u_1 - u_2 + u_3\bigr),
\qquad
X^2 = u_2 - 1 = \tfrac12\bigl(-u_1 + u_2 + u_3\bigr).
\]
والوحيدات الحدّ تقع في \omterm{def:b1:vspaces:span}{الفضاء المولَّد}، ومنه فكل شيء كذلك: ومنه فالعائلة $\mathcal{F}$
\omterm{def:b1:vspaces:free}{أساس}. وكعلاوة، يعطي تجميع المعروضات الثلاث
\omterm{prop:b1:vspaces:coordinates}{إحداثيات} أيّ $P = \alpha + \beta X + \gamma X^2$:
\[
P = \frac{\alpha + \beta - \gamma}{2}\,u_1
+ \frac{\alpha - \beta + \gamma}{2}\,u_2
+ \frac{-\alpha + \beta + \gamma}{2}\,u_3 .
\]
(وللتحقق من السلامة مع $P = X$: \omterm{prop:b1:vspaces:coordinates}{الإحداثيات} $\bigl(\frac12,
-\frac12, \frac12\bigr)$، كما وُجد أعلاه.) ودرسان: عادةً ما يخفي التناظر
في العائلة تركيبةً مختصرة؛ وما إن يتوفر البُعد
(\cref{ch:b1:findim})، يأتي نصف التوليد كله
من هذا العمل بالمجّان --- فثلاث متجهات \omterm{def:b1:vspaces:free}{حرة} في فضاء
بُعده $3$ تكوّن \omterm{def:b1:vspaces:free}{أساسًا} دائمًا.
\end{example}

\begin{proposition}[محكات حرية نافعة]\label{prop:b1:vspaces:freecriteria}
\begin{enumerate}
  \item كل عائلة من \emph{كثيرات حدود غير معدومة ذات درجات متمايزة
        مثنى مثنى} حرةٌ.
  \item إضافة متجهة إلى \omterm{def:b1:vspaces:free}{عائلة حرة} تبقيها \omterm{def:b1:vspaces:free}{حرة} إذا وفقط
        إذا كانت المتجهة خارج \omterm{def:b1:vspaces:span}{الفضاء المولَّد} بالعائلة.
  \item كل عائلة جزئية من \omterm{def:b1:vspaces:free}{عائلة حرة} حرةٌ؛ وكل عائلة
        تحتوي عائلة مولِّدة مولِّدةٌ.
\end{enumerate}
\end{proposition}

\begin{proof}
(1) في تركيبة معدومة، انظر في أكبر درجة حاضرة: فعلى
معاملها أن ينعدم (لأن لا شيء يختصر تلك الدرجة)، وتتالَ
نزولًا.

(2) إذا كان $x \in \operatorname{Vect}(x_1, \dots, x_p)$، فإن العلاقة $x
- \sum\lambda_i x_i = 0$ غير بديهية. وبالعكس، على تركيبة معدومة
غير بديهية للعائلة $(x_1, \dots, x_p, x)$ أن تتضمن $x$ بمعامل
غير معدوم (وإلا ناقضت حرية العائلة
الصغرى)، والحلّ من أجل $x$ يضعه في \omterm{def:b1:vspaces:span}{الفضاء المولَّد}.

(3) العائلة الجزئية: تركيبة معدومة للعائلة الجزئية هي تركيبة
للعائلة كلها بجعل المعاملات الناقصة $0$؛ وتقتل حرية
العائلة الكبرى كلَّها. والعائلة الفائقة: كل متجهة
من $E$ هي أصلًا تركيبة من الجزء المولِّد؛ فأعطِ
المتجهات الإضافية المعامل $0$.
\end{proof}

\begin{example}[مبدأ السلّم]
\label{ex:b1:vspaces:staircase}
لتكن $P_0, P_1, \dots, P_n \in K_n[X]$ مع $\deg P_k = k$ من أجل كل
$k$ (أي «سلّم» درجات). عندئذ تكون $(P_0, \dots, P_n)$
\omterm{def:b1:vspaces:free}{أساسًا} للمقدار $K_n[X]$. والحرية هي
\cref{prop:b1:vspaces:freecriteria}~(1). وأمّا خاصية
التوليد، فحاجج بنزول منته على الدرجة: لتكن $Q \in
K_n[X]$ مع $Q \neq 0$، من الدرجة $d$، ومعاملها المهيمن $a$،
وليكن $b \neq 0$ المعامل المهيمن للمقدار $P_d$. عندئذ يكون $Q -
\frac ab P_d$ من الدرجة $< d$ (لأن الحدود العليا تتلاشى)؛ وباستبدال
هذا الفرق بالمقدار $Q$ والتكرار، يُبلغ بعد $n + 1$ خطوة على الأكثر
كثيرُ الحدود المعدوم، ويعبّر فكّ الطروح عن $Q$
تركيبةً من المقادير $P_k$. وقد لُقي سلّمان أصلًا:
القوى المنسحبة $\bigl((X-a)^k\bigr)_{0 \leq k
\leq n}$ (\cref{exo:b1:vspaces:4})، وجداءات نيوتن
$\bigl((X - x_0)(X - x_1)\cdots(X - x_{k-1})\bigr)_{0 \leq k \leq
n}$، المشغَّلة في مسألة نهاية الأسبوع.
\end{example}

\begin{example}[الحرية في فضاءات الدوال]\label{ex:b1:vspaces:functionfree}
في $\mathcal{F}(\R,\R)$، تكون العائلة $(\eu^{a_1 x}, \dots, \eu^{a_p
x})$ مع $a_1 < \dots < a_p$ حرةً: اقسم تركيبة معدومة على
$\eu^{a_p x}$ ودع $x \to +\infty$؛ فيموت المعامل الأخير،
ويتتالى المرء نزولًا (\cref{exo:b1:vspaces:8} يفصّل هذا و
تنويعاته). فحرية الدوال يُبرهن عليها عن طريق \emph{التقويم}:
عند نقاط محسنة الاختيار، أو عند اللانهاية، أو بعد الاشتقاق.
\end{example}

\begin{example}[علاقة مستترة تقلّص فضاءً مولَّدًا]
\label{ex:b1:vspaces:hiddenrelation}
في $\mathcal{F}(\R, \R)$، ما هو
$\operatorname{Vect}(1,\ \cos^2,\ \sin^2)$؟ المتطابقة
$\cos^2 + \sin^2 = 1$ تركيبةٌ معدومة غير بديهية
\[
1\cdot\mathbf{1} + (-1)\cos^2 + (-1)\sin^2 = 0 :
\]
ومنه فالعائلة مرتبطة، \omterm{def:b1:vspaces:span}{والفضاء المولَّد} مولَّد أصلًا
بالعائلة $(1, \cos^2)$ وحدها (لأن $\sin^2 = 1 - \cos^2$). وتلك العائلة الأصغر
\omterm{def:b1:vspaces:free}{حرة}: إذ يعطي $a + b\cos^2 x = 0$ من أجل كل $x$، عند $x = 0$ و
$x = \frac\pi2$: $a + b = 0$ و $a = 0$. ومنه \omterm{def:b1:vspaces:span}{فالفضاء المولَّد}
\emph{مستوٍ} داخل فضاء الدوال --- وهو يحتوي كذلك
$\cos 2x = 2\cos^2 x - 1$: فالعائلات التي تبدو خطية من
الدوال المثلثية تنهار روتينيًا تحت المتطابقات،
ولهذا يجب \emph{البرهان} على الحرية، لا افتراضها من
طول القائمة.
\end{example}

\begin{remark}[مزالق شائعة]
\label{rem:b1:vspaces:pitfalls}
أربعة فخاخ كلاسيكية.
\emph{المثنى مثنى لا يكفي}: ففي $\R^2$، المتجهات $(1,0)$
و $(0,1)$ و $(1,1)$ غير متناسبة مثنى مثنى، ومع ذلك مرتبطة ---
فالحرية خاصية للعائلة \emph{كلها}، تُختبر بتركيبة
شاملة واحدة، لا اثنين اثنين أبدًا.
\emph{المتجهة المعدومة تسمّم كل شيء}: فكل عائلة تحتوي
$0$ مرتبطة (لأن $1\cdot 0 = 0$ علاقة غير بديهية)، مهما كانت
المتجهات الأخرى بريئة.
\emph{الاتحاد ليس المجموع}: فالمقدار $F \cup G$ ليس \omterm{def:b1:vspaces:subspace}{فضاءً جزئيًا} في الغالب الأعمّ
(\cref{def:b1:vspaces:subspace})؛ وأصغر \omterm{def:b1:vspaces:subspace}{فضاء جزئي}
يحتوي كليهما هو $F + G$، وهو عادةً أكبر بكثير من الاتحاد
--- ففي $\R^2$، لمستقيمين متمايزين اتحادٌ متقاطع، ومجموعٌ هو
المستوي كله.
\emph{المباشر يقتضي تقاطعًا بديهيًا لا انفصالًا}:
فالفضاءان الجزئيان ليسا منفصلين أبدًا (لأن كليهما يحتوي $0$)؛ والشرط
الصحيح هو $F \cap G = \{0\}$، ويجب \emph{البرهان} عليه،
لا قراءته من رسم --- قارن
\cref{ex:b1:vspaces:manysupplements}، حيث يفي كثير من المقادير $G$ المختلفة بالغرض.
\emph{الحرية تتعلق بالسلالم}: فالزوج $(1, \iu)$
حرّ في $\C$ منظورًا إليه فضاءً متجهيًا على $\R$، لكنه مرتبط في $\C$
منظورًا إليه فضاءً متجهيًا على $\C$ ($\iu\cdot 1 + (-1)\cdot\iu = 0$).
فاعلم دائمًا أيّ \omterm{def:b1:structures:field}{حقل} يفعل قبل أن تعلن عائلةً \omterm{def:b1:vspaces:free}{حرة}
--- وتحوّل مسألة نهاية الأسبوع في \cref{ch:b1:findim} هذه
الحساسية بالضبط إلى براهين صمم.
\end{remark}

\begin{remark}[إلى أين تذهب هذه اللغة]
\label{rem:b1:vspaces:whereused}
كل ما بعد هذا الفصل يتكلم اللغة المهيّأة هنا.
فالفصل \cref{ch:b1:findim} يعدّ متجهات \omterm{def:b1:vspaces:free}{الأساس} ويحوّل «\omterm{def:b1:vspaces:free}{الحرة}» و
«المولِّدة» إلى متراجحات على عدد صحيح واحد، هو
البُعد. و \cref{ch:b1:linmaps} يدرس التطبيقات المتوافقة مع
العمليتين؛ وتصير المجاميع المباشرة مساقطَ هناك.
و \cref{ch:b1:matrices} يرمّز المتجهات بإحداثياتها في
\omterm{def:b1:vspaces:free}{أساس} --- و \cref{prop:b1:vspaces:coordinates} هو الرخصة لهذا
الترميز --- ويضيف \cref{ch:b1:euclid} الأطوال والزوايا
فوق البنية الخطية. وفي مجلّد السنة~الثانية تجري البديهيات
نفسها، حرفيًا، على حقول كيفية وفي بُعد غير
منته؛ فلم يستعمل أيّ شيء في هذا الفصل التناهي في أيّ موضع.
\end{remark}

\begin{remark}[ثلاثة خيوط تُتبع عبر الكتاب 3]
\label{rem:b1:vspaces:threads}
راقب ثلاث أفكار محدَّدة من هذا الفصل تنمو.
\emph{مبدأ السلّم}
(\cref{ex:b1:vspaces:staircase}) يعود أساسَ نيوتن
في مسألة نهاية الأسبوع في هذا الفصل، \omterm{def:b1:vspaces:free}{والأساس} الثنائي
$(B_k)$ هناك، وحيلةَ المتناوب الكثيرالحدود في
مسألة نهاية الأسبوع في \cref{ch:b1:det}: أي مبرهنة مساعدة واحدة وثلاثة
أرباح بلا محددات.
\emph{والتقويم محكَّ حرية}
(\cref{ex:b1:vspaces:functionfree}) يصير تشاكل المقايسة التقابلي
في \cref{ch:b1:linmaps}، ثم محك فاندرموند
في \cref{ch:b1:det}، ثم اختبار غرام في
\cref{ch:b1:euclid}: أي المنعكس نفسه، محدَّدًا ثلاث مرات.
\emph{والمجاميع المباشرة} (\cref{def:b1:vspaces:sum}) تصير
مساقط في \cref{ch:b1:linmaps}، وتفكيكات متعامدة $E =
F \oplus F^\perp$ في \cref{ch:b1:euclid}، والتفكيك إلى مفسَّر زائد باقٍ
في المربعات الصغرى في مسألة نهاية الأسبوع في \cref{ch:b1:multivar}. وقليل جدًا من هذا
الكتاب ليس، في أساسه، واحدة من هذه الأفكار الثلاث في ثياب
جديدة.
\end{remark}

\section{تمارين}

\begin{exercise}[$\star$]\label{exo:b1:vspaces:1}
أيّ ممّا يأتي فضاءات جزئية؟
\begin{enumerate}
  \item $\{(x, y, z) \in \R^3 : x + 2y - z = 0\}$؛
  \item $\{(x, y, z) \in \R^3 : x + 2y - z = 1\}$؛
  \item $\{(x, y) \in \R^2 : xy \geq 0\}$؛
  \item $\{P \in \R[X] : P(1) = 0\}$؛
  \item $\{f \in \mathcal{F}(\R,\R) : f \text{ محدودة}\}$.
\end{enumerate}
\end{exercise}

\begin{exercise}[$\star$]\label{exo:b1:vspaces:2}
في $\R^3$، هل $(1, 2, 1)$ في
$\operatorname{Vect}\bigl((1,0,1),\, (1,1,0)\bigr)$؟ وماذا عن $(2, 1,
1)$؟ وصف $\operatorname{Vect}\bigl((1,0,1),(1,1,0)\bigr)$ بمعادلة.
\end{exercise}

\begin{exercise}[$\star$]\label{exo:b1:vspaces:3}
افصل في الحرية في $\R^3$: $\;\bigl((1,1,0), (1,0,1),
(0,1,1)\bigr)$؛ $\;\bigl((1,2,3), (2,4,6)\bigr)$؛
$\;\bigl((1,0,0), (1,1,0), (1,1,1), (0,1,1)\bigr)$.
\end{exercise}

\begin{exercise}[$\star$]\label{exo:b1:vspaces:4}
برهن على أن $(1, X - 1, (X-1)^2, (X-1)^3)$ \omterm{def:b1:vspaces:free}{أساس} للمقدار $\R_3[X]$،
وأعطِ \omterm{prop:b1:vspaces:coordinates}{إحداثيات} $X^3$ فيه. \emph{(تايلور عند $1$!)}
\end{exercise}

\begin{exercise}[$\star\star$]\label{exo:b1:vspaces:5}
في $\R^4$، لتكن $F = \{(x,y,z,t) : x = y = z\}$ و $G = \{(x,y,z,t)
: x = t = 0\}$. برهن على أن $F \oplus G = \R^4$، وفكّك
$(1,2,3,4)$ تبعًا لذلك.
\end{exercise}

\begin{exercise}[$\star\star$]\label{exo:b1:vspaces:6}
في فضاء المتتاليات، لتكن $F$ \omterm{def:b1:logic:sets}{مجموعة} المتتاليات
المتقاربة و $G = \operatorname{Vect}(u)$ حيث $u_n = (-1)^n$.
برهن على أن $F \cap G = \{0\}$. وهل $F + G$ هو فضاء المتتاليات
كله؟
\end{exercise}

\begin{exercise}[$\star\star$]\label{exo:b1:vspaces:7}
لتكن $F, G, H$ فضاءات جزئية من $E$. برهن على أن
\[
F \cap (G + (F \cap H)) = (F \cap G) + (F \cap H),
\]
وبيّن بمثال في $\R^2$ أن التوزيعية غير المقيَّدة
$F \cap (G + H) = (F\cap G) + (F \cap H)$ تفشل.
\end{exercise}

\begin{exercise}[$\star\star\star$]\label{exo:b1:vspaces:8}
برهن على أن العائلات الآتية من $\mathcal{F}(\R, \R)$ \omterm{def:b1:vspaces:free}{حرة}:
\begin{enumerate}
  \item $(\eu^{a_1 x}, \dots, \eu^{a_p x})$ من أجل $a_1 < \dots < a_p$؛
  \item $(\cos x, \sin x, \cos 2x, \sin 2x)$؛
  \item $(x \mapsto \abs{x - a_1}, \dots, x \mapsto \abs{x - a_p})$
        من أجل $a_i$ متمايزة \emph{(فالقابلية للاشتقاق تفشل عند نقطة
        واحدة بالضبط لكل دالة)}.
\end{enumerate}
\end{exercise}

\begin{exercise}[$\star\star\star$]\label{exo:b1:vspaces:9}
ليكن $E$ فضاءً متجهيًا على $K$ ولتكن $F, G, H$ فضاءات جزئية مع $F + G =
F + H$ و $F \cap G = F \cap H$ و $G \subseteq H$. برهن على $G = H$.
وأعطِ مثالًا مضادًا دون الفرض $G \subseteq H$.
\end{exercise}

\begin{exercise}[$\star\star$]\label{exo:b1:vspaces:10}
في $\R[X]$، لتكن $\mathcal P$ \omterm{def:b1:logic:sets}{مجموعة} كثيرات الحدود \emph{الزوجية}
($P(-X) = P(X)$) و $\mathcal I$ \omterm{def:b1:logic:sets}{مجموعة} \emph{الفردية}
($P(-X) = -P(X)$). برهن على أن $\R[X] = \mathcal P \oplus \mathcal
I$، وبيّن أن $\mathcal P =
\operatorname{Vect}(1, X^2, X^4, \dots)$، أي أن كثيرات الحدود
الزوجية هي بالضبط كثيرات الحدود في $X^2$.
\end{exercise}

\begin{exercise}[$\star\star$]\label{exo:b1:vspaces:11}
لتكن $(x_1, x_2, x_3)$ \omterm{def:b1:vspaces:free}{عائلة حرة} من \omterm{def:b1:vspaces:def}{فضاء متجهي} حقيقي $E$.
برهن على أن $(x_1 + x_2,\ x_2 + x_3,\ x_3 + x_1)$ \omterm{def:b1:vspaces:free}{حرة}. وهل تكون
العائلة النظيرة ذات المتجهات الأربع $(x_1 + x_2,\ x_2 + x_3,\ x_3 +
x_4,\ x_4 + x_1)$ \omterm{def:b1:vspaces:free}{حرة} عندما تكون $(x_1, x_2, x_3, x_4)$ كذلك؟
\end{exercise}

\begin{exercise}[$\star\star\star$]\label{exo:b1:vspaces:12}
ليكن $E$ فضاءً متجهيًا على $\R$ (أو $\C$) ولتكن $F_1, \dots,
F_k$ فضاءات جزئية \emph{قطعية} من $E$ (أي كل $F_i \neq E$).
\begin{enumerate}
  \item عالج الحالة $k = 2$ مباشرة: إذا كان $F_1 \not\subseteq
        F_2$ و $F_2 \not\subseteq F_1$، فاختر $x \in F_1
        \setminus F_2$ و $y \in F_2 \setminus F_1$ وحدّد موضع $x
        + y$.
  \item برهن عمومًا على أن $E \neq F_1 \cup \dots \cup F_k$: أي إن
        \omterm{def:b1:vspaces:def}{الفضاء المتجهي} على \omterm{def:b1:structures:field}{حقل} غير منته ليس أبدًا اتحادًا منتهيًا
        لفضاءات جزئية قطعية. \emph{(خذ $k$ أصغريًا، واختر
        $x \in F_1$ خارج بقية $F_i$، واختر $y \notin F_1$،
        وتتبّع المستقيم $t \mapsto y + tx$.)}
\end{enumerate}
\end{exercise}

\section{مسألة: المقايسة، ثلاثة أسس لفضاء واحد}

\begin{problem}\label{pb:b1:vspaces:1}
ثبّت $n + 1$ نقطة \emph{متمايزة} $x_0, x_1, \dots, x_n$ من $\R$.
وتعيد هذه المسألة النظر في \omterm{thm:b1:poly:lagrange}{مقايسة لاغرانج}
(\cref{thm:b1:poly:lagrange}) بعينَي هذا الفصل: فالفضاء
$\R_n[X]$ يحمل ثلاثة أسس طبيعية --- \omterm{def:b1:vspaces:free}{أساس} لاغرانج،
\omterm{def:b1:vspaces:free}{وأساس} نيوتن، و(من أجل النقاط متساوية التباعد) \omterm{def:b1:vspaces:free}{الأساس} الثنائي ---
ويجعل كل \omterm{def:b1:vspaces:free}{أساس} سؤالًا واحدًا سهلًا. وينتهي الطريق إلى مبرهنة
حسابية حقيقية: تمييز بوليا لكثيرات الحدود
التي ترسل $\Z$ داخل $\Z$.\index{مقايسة لاغرانج}
\index{فروق مقسومة}\index{كثير حدود صحيح القيم}

\medskip
\noindent\textbf{الجزء 1 --- \omterm{def:b1:vspaces:free}{أساس} لاغرانج.} من أجل $0 \leq i
\leq n$ ضع
\[
L_i \;=\; \prod_{j \neq i} \frac{X - x_j}{x_i - x_j} \;\in\;
\R_n[X].
\]
\begin{enumerate}
  \item تحقق من أن $\deg L_i = n$ ومن أن $L_i(x_j) = 1$ إذا كان $j =
        i$، و $0$ إذا كان $j \neq i$.
  \item برهن على أن العائلة $(L_0, \dots, L_n)$ \omterm{def:b1:vspaces:free}{حرة}.
  \item برهن على أنه من أجل كل $P \in \R_n[X]$،
        \[
        P \;=\; \sum_{i=0}^{n} P(x_i)\, L_i ,
        \]
        واستنتج أن $(L_0, \dots, L_n)$ \omterm{def:b1:vspaces:free}{أساس} للمقدار
        $\R_n[X]$. \emph{(انظر في فرق الطرفين
        وعُدّ جذوره، \cref{cor:b1:poly:nroots}.)}
  \item استنتج مبرهنة المقايسة: من أجل أيّ قيم $y_0,
        \dots, y_n \in \R$ يوجد $P \in
        \R_n[X]$ \emph{وحيد} يحقق $P(x_i) = y_i$ من أجل كل $i$. وفي \omterm{def:b1:vspaces:free}{أساس}
        لاغرانج، ما \omterm{prop:b1:vspaces:coordinates}{إحداثيات} \omterm{def:b1:poly:def}{كثير حدود} $P$؟
  \item برهن على المتطابقتين
        \[
        \sum_{i=0}^{n} L_i = 1
        \qquad\text{و، من أجل } 0 \leq k \leq n,\qquad
        \sum_{i=0}^{n} x_i^{k}\, L_i = X^{k} .
        \]
\end{enumerate}

\medskip
\noindent\textbf{الجزء 2 --- \omterm{def:b1:vspaces:free}{أساس} نيوتن \omterm{pb:b1:vspaces:1}{والفروق
المقسومة}.} ضع $N_0 = 1$ و $N_k = (X - x_0)(X - x_1) \cdots
(X - x_{k-1})$ من أجل $1 \leq k \leq n$. ومن أجل دالة $f$ معرَّفة
عند العقد، عرّف \emph{\omterm{pb:b1:vspaces:1}{الفروق المقسومة}} بالمقدار
$f[x_i] = f(x_i)$ و
\[
f[x_i, \dots, x_{i+k}] \;=\;
\frac{f[x_{i+1}, \dots, x_{i+k}] - f[x_i, \dots, x_{i+k-1}]}
{x_{i+k} - x_i} .
\]
\begin{enumerate}[resume]
  \item برهن على أن $(N_0, N_1, \dots, N_n)$ \omterm{def:b1:vspaces:free}{أساس} للمقدار
        $\R_n[X]$.
  \item احسب $f[x_0, x_1]$ و $f[x_0, x_1, x_2]$ بدلالة
        قيم $f$، ثم احسب كل \omterm{pb:b1:vspaces:1}{الفروق المقسومة} للمقدار
        $f(x) = x^2$ عند ثلاث عقد كيفية.
  \item (مبرهنة أيتكن المساعدة) ليقايس $R$ الدالةَ $f$ عند $x_0, \dots,
        x_{n-1}$ وليقايس $Q$ الدالةَ $f$ عند $x_1, \dots, x_n$،
        وكلاهما من الدرجة $\leq n - 1$. برهن على أن
        \[
        S \;=\; \frac{(X - x_0)\,Q - (X - x_n)\,R}{x_n - x_0}
        \]
        يقايس $f$ عند $x_0, x_1, \dots, x_n$.
  \item استنتج، بالاستقراء على عدد العقد، أن
        معامل $X^{k}$ في مقايِس $f$ عند $x_0,
        \dots, x_k$ هو بالضبط $f[x_0, \dots, x_k]$.
  \item برهن على \emph{صيغة نيوتن للمقايسة}: مقايِس $f$
        عند $x_0, \dots, x_n$ هو
        \[
        P \;=\; \sum_{k=0}^{n} f[x_0, \dots, x_k]\, N_k ,
        \]
        واستخرج الصيغة \omterm{def:b1:topology:closed}{المغلقة}
        \[
        f[x_0, \dots, x_k] \;=\; \sum_{i=0}^{k}
        \frac{f(x_i)}{\prod_{j \neq i,\, j \leq k} (x_i - x_j)} ,
        \]
        التي تبيّن أن $f[x_0, \dots, x_k]$ لا يتعلق
        بترتيب العقد.
\end{enumerate}

\medskip
\noindent\textbf{الجزء 3 --- العقد متساوية التباعد: مؤثر
الفرق.} من الآن فصاعدًا تكون العقد $0, 1, 2, \dots$ و، من أجل
\omterm{def:b1:poly:def}{كثير حدود} $P$، نضع
\[
\Delta P(X) = P(X + 1) - P(X),
\qquad
B_k = \frac{X(X-1)\cdots(X-k+1)}{k!} \quad (B_0 = 1).
\]
\begin{enumerate}[resume]
  \item بيّن أنه إذا كان $\deg P = m \geq 1$ بمعامل مهيمن
        $a$، فإن $\deg \Delta P = m - 1$ بمعامل مهيمن
        $m\,a$، وأن $\Delta$ يقتل الثوابت.
  \item بيّن أن $(B_0, B_1, \dots, B_n)$ \omterm{def:b1:vspaces:free}{أساس} للمقدار
        $\R_n[X]$ وأن $\Delta B_k = B_{k-1}$ من أجل $k \geq 1$.
  \item (صيغة نيوتن للفروق الأمامية) برهن على أن كل $P
        \in \R_n[X]$ يحقق
        \[
        P \;=\; \sum_{k=0}^{n} \bigl(\Delta^{k} P\bigr)(0)\, B_k .
        \]
  \item برهن على أنه من أجل كل $k \geq 0$،
        \[
        \bigl(\Delta^{k} P\bigr)(0) \;=\;
        \sum_{j=0}^{k} (-1)^{k-j} \binom{k}{j} P(j) .
        \]
  \item بيّن أنه إذا كان $\deg P = n$ بمعامل مهيمن $a_n$،
        فإن $\Delta^{n} P$ هو الثابت $n!\,a_n$ و
        $\Delta^{n+1} P = 0$.
\end{enumerate}

\medskip
\noindent\textbf{الجزء 4 --- كثيرات الحدود صحيحة القيم.} يكون
\omterm{def:b1:poly:def}{كثير الحدود} $P \in \R[X]$ \emph{صحيح القيم} إذا كان $P(m) \in
\Z$ من أجل كل $m \in \Z$.
\begin{enumerate}[resume]
  \item برهن على أن كل $B_k$ صحيح القيم. \emph{(عالج $m
        \geq k$ و $0 \leq m < k$ و $m < 0$ على حدة؛ ومن أجل $m =
        -q < 0$، بيّن $B_k(-q) = (-1)^k \binom{q + k - 1}{k}$.)}
  \item برهن على \emph{تمييز بوليا}: يكون $P \in \R_n[X]$
        صحيح القيم إذا وفقط إذا كانت إحداثياته في
        \omterm{def:b1:vspaces:free}{الأساس} $(B_0, \dots, B_n)$ أعدادًا صحيحة.
  \item استنتج: إذا أخذ $P \in \R_n[X]$ قيمًا صحيحة عند $n + 1$
        عددًا صحيحًا \emph{متتاليًا} $a, a+1, \dots, a+n$، فإن $P$
        صحيح القيم. \emph{(بالانسحاب: طبّق الدراسة على $Q(X)
        = P(X + a)$.)}
  \item استنتج من السؤال 16 أن جداء $k$ عددًا صحيحًا
        متتاليًا يقبل القسمة على $k!$ دائمًا.
  \item لتكن $P = \dfrac{X(X+1)(2X+1)}{6}$. احسب جدول نيوتن لها
        عند $0, 1, 2, 3$، واكتب $P$ في \omterm{def:b1:vspaces:free}{الأساس} $(B_k)$، و
        استنتج أن $P$ صحيح القيم وإن لم يكن أيّ من معاملاتها
        الوحيدة الحدّ عددًا صحيحًا. وتحقق من $\Delta P =
        (X+1)^2$ واستنتج $P(m) = 1^2 + 2^2 + \dots + m^2$ من أجل
        $m \in \N$.
\end{enumerate}

\medskip
\noindent\textbf{الجزء 5 --- الأرباح.}
\begin{enumerate}[resume]
  \item ليقايس $P \in \R_n[X]$ القيمَ $2^i$ عند $i =
        0, 1, \dots, n$. بيّن أن $P = B_0 + B_1 + \dots + B_n$
        وأن $P(n + 1) = 2^{n+1} - 1$: أي إن «نمط
        المضاعفة» ينكسر دائمًا عند النقطة التالية مباشرة.
  \item (الدالة الأصلية المتقطّعة) برهن على أنه من أجل كل عددين صحيحين $m
        \geq 1$ و $k \geq 0$،
        \[
        \sum_{j=0}^{m-1} B_k(j) \;=\; B_{k+1}(m),
        \]
        أي متطابقة عصا الهوكي $\sum_{j=k}^{m-1}
        \binom{j}{k} = \binom{m}{k+1}$.
  \item انشر $X^2$ و $X^3$ في \omterm{def:b1:vspaces:free}{الأساس} $(B_k)$ واستنتج
        صيغتين مغلقتين من أجل $\sum_{j=0}^{m-1} j^2$ و
        $\sum_{j=0}^{m-1} j^3$؛ واستعد متطابقة نيقوماخوس
        $1^3 + \dots + m^3 = (1 + \dots + m)^2$.
  \item خذ $n = 2$ والعقد $0, 1, 2$. اكتب \omterm{prop:b1:vspaces:coordinates}{إحداثيات}
        $X^2$ في أسس هذه المسألة الثلاثة: \omterm{def:b1:vspaces:free}{الأساس} الوحيد الحدّ،
        \omterm{def:b1:vspaces:free}{وأساس} لاغرانج، \omterm{def:b1:vspaces:free}{وأساس} نيوتن. وتحقق من
        الأجوبة الثلاثة إزاء السؤالين 4 و 9.
  \item توليفة. في أربع جمل: أيّ مفهوم من مفاهيم الفضاءات المتجهية
        يجعل السؤال 4 تلقائيًا؛ ولماذا يحسب \omterm{def:b1:vspaces:free}{أساس} نيوتن
        \omterm{prop:b1:vspaces:coordinates}{الإحداثيات} \emph{تراجعيًا} بينما يقرؤها \omterm{def:b1:vspaces:free}{أساس} لاغرانج
        \emph{فورًا}؛ وأيّ محك حرية يتقاسمه
        الأساسان؛ وبأيّ معنى دقيق تقول مبرهنة بوليا إن
        صحيحية \omterm{def:b1:poly:def}{كثير حدود} خاصيةٌ لإحداثياته
        \emph{في \omterm{def:b1:vspaces:free}{الأساس} الصحيح}.
\end{enumerate}
\end{problem}
